%!TEX root = ../template.tex
%%%%%%%%%%%%%%%%%%%%%%%%%%%%%%%%%%%%%%%%%%%%%%%%%%%%%%%%%%%%%%%%%%%%
%% abstract-en.tex
%% NOVA thesis document file
%%
%% Abstract in English
%%%%%%%%%%%%%%%%%%%%%%%%%%%%%%%%%%%%%%%%%%%%%%%%%%%%%%%%%%%%%%%%%%%%

\typeout{NT FILE abstract-en.tex}%

{
  The management of internships, dissertations, and engineering projects at the Department of Computer Science (DI) of NOVA School of Science and Technology (FCT)
  currently relies on a legacy platform that is technically obsolete and limited to undergraduate cycles.
  Master's degree processes are supported by spreadsheets and informal email exchanges, leading to fragmented data, audit difficulties, and high administrative burden
  for coordinators, professors, companies, and the secretariat.
  Simultaneously, the existing platform suffers from UX problems, email deliverability failures, and security vulnerabilities, making incremental evolution unfeasible.
}

{
  This problem is relevant and challenging because internship placement is a critical process for students' academic paths and for the department's relationship
  with companies.
  The solution must reconcile complex functional requirements (automated seriation, multiple user profiles, distinct workflows for different degrees) with demanding
  non-functional requirements (security, auditability, interoperability, and usability) in an institutional context where long-term maintainability is as important as development speed.
}

{
  This dissertation proposes the design and implementation of a new unified platform for managing proposals and candidacies, based on a layered architecture: a React frontend,
  a backend initially supported by Supabase and subsequently evolved to Spring Boot, and a PostgreSQL database. 
  The work follows an iterative methodology: requirements gathering, process (BPMN) and data (ERD) modeling, use case definition, and accelerated prototyping through Vibe Coding tools, 
  followed by refactoring and architectural hardening.
}

{
  As main results, a functional prototype is expected to automate critical tasks (seriation, notifications, protocols), improve user experience and data quality, and provide a critical
  analysis of the user of Vibe Coding tools in institutional contexts.
}


\keywords{
  Internship and dissertation management \and
  Web Platforms \and
  Software Architectures \and
  Vibe Coding \and
  Full-Stack Development
}

%Regardless of the language in which the dissertation is written, usually there are at least two abstracts: one abstract in the same language as the main text, and another abstract in some other language.

%The abstracts' order varies with the school.  If your school has specific regulations concerning the abstracts' order, the \gls{novathesis} (\LaTeX) template will respect them.  Otherwise, the default rule in the \gls{novathesis} template is to have in first place the abstract in \emph{the same language as main text}, and then the abstract in \emph{the other language}. For example, if the dissertation is written in Portuguese, the abstracts' order will be first Portuguese and then English, followed by the main text in Portuguese. If the dissertation is written in English, the abstracts' order will be first English and then Portuguese, followed by the main text in English.
%
%However, this order can be customized by adding one of the following to the file \verb+5_packages.tex+.

%\begin{verbatim}
%    \ntsetup{abstractorder={<LANG_1>,...,<LANG_N>}}
%    \ntsetup{abstractorder={<MAIN_LANG>={<LANG_1>,...,<LANG_N>}}}
%\end{verbatim}

%For example, for a main document written in German with abstracts written in German, English and Italian (by this order) use:
%\begin{verbatim}
%    \ntsetup{abstractorder={de={de,en,it}}}
%\end{verbatim}

%Concerning its contents, the abstracts should not exceed one page and may answer the following questions (it is essential to adapt to the usual practices of your scientific area):

%\begin{enumerate}
%  \item What is the problem?
%  \item Why is this problem interesting/challenging?
%  \item What is the proposed approach/solution/contribution?
%  \item What results (implications/consequences) from the solution?
%\end{enumerate}

% Abstract keywords
%\keywords{
%  One keyword \and
%  Another keyword \and
%  Yet another keyword \and
%  One keyword more \and
%  The last keyword
%}
